\scanpage{2}
\mspage{1}
\textcircled{1} Soit $\overline{\mathbf Q}$ la clôture algébrique de $\mathbf Q$ dans $\mathbf C$.
Si $X=\overline X$ est un schéma sur $\overline{\mathbf Q}$, on considère le
groupe $\mathcal G_{\overline X}$ des automorphismes de $\overline X$ en tant
que schéma absolu. On a une suite
\[
(1)\qquad 1\longrightarrow G_{\overline X}\longrightarrow\mathcal G_{\overline X}\longrightarrow\Gamma
\]
% The upper unresolved cluster and the lower unreadable span are distinct.
\margin{\begin{tabular}{@{}l@{}}
Question : \unsure{\tikz[baseline=-.4ex,x=.6em,y=.6em]{
  \draw (.10,1.20) .. controls (.45,1.75) and (.85,1.70) .. (.70,1.20)
    .. controls (.45,.70) and (.20,.15) .. (-.05,-1.10);
  \draw (.65,1.25) .. controls (1.10,1.65) and (1.20,1.25) .. (.80,.98)
    --(.55,-.45) .. controls (.60,-.10) and (1.12,-.15) .. (1.15,-.50);
  \draw (-.10,.85)--(.80,1.08);
  \draw (-.15,.60) .. controls (.45,.70) and (.65,.55) .. (1.00,.80);
}}\\
\unsure{comment|comme} \illegible{1} ?
\end{tabular}}
où
\[
(2)\qquad\left\{\begin{aligned}
\Gamma&=\operatorname{Gal}(\overline{\mathbf Q}/\mathbf Q)\\
G_{\overline X}&:\operatorname{Aut}(\overline X/\overline{\mathbf Q})
\end{aligned}\right.
\]
% The small subset-shaped curve is detached below the arrow shaft.
\margin{\diagram{\begin{tikzcd}[ampersand replacement=\&,column sep=small]
\overline{\mathbf Q}\arrow[d,phantom,"\cup" description]
\&\mathrel{\tikz[baseline=-.45ex,x=1em,y=1em]{
  \draw[->] (0,0)--(1.8,0);
  \draw (.30,-.12) .. controls (.02,-.12) and (.02,-.35) .. (.30,-.35);
}}
\&\Gamma(X_{\mathrm{red}},\underline{\mathcal O}_X)\\
\mathbf Q\&\&
\end{tikzcd}}
Il faut supposer p.ex. que $\overline{\mathbf Q}$ est la
clôt. int. de $\mathbf Q$ ds $\Gamma_X$}
(3) Les sous-extensions $K$ de $\overline{\mathbf Q}/\mathbf Q$ correspondent
aux sous-groupes fermés $\Gamma_K$ de $\Gamma$ (d'indice
fini ssi $K$ fini sur $\mathbf Q$). À une restriction
du groupe de base $\overline{\mathbf Q}$ de $\overline X$ \unsure{:|à} $K$ correspond
un relèvement
% The upper connector has a shaft and an upper terminal cap, not a complete two-wing head.
\diagram{\begin{tikzcd}[ampersand replacement=\&,column sep=small]
(4)\quad\mathcal G_{\overline X}
\&\mathrel{\tikz[baseline=-.45ex,x=1em,y=1em]{
  \draw (0,0)--(1.8,0);
  \fill (1.8,0)--(1.60,.12)--(1.60,0)--cycle;
}}
\&\Gamma\arrow[d,phantom,"\cup" description]\\
\&\&\Gamma_K\arrow[ull]
\end{tikzcd}\,,}
\margin{grain de sel d'effectivité de descente
\unsure{et l'hypothèse}
\unsure{qu. proj.|gr. proj.}}
qui le définit si $\overline X$ est
\tikz[baseline=(w.base)]{
  \node[inner sep=0pt] (w) at (0,0) {réduit};
  \node[inner sep=0pt] at (.48,.20) {?};
  \draw (.08,.06) ellipse (.73cm and .32cm);
}
et \uline{loc de t f.}
sur $\overline{\mathbf Q}$. Cela montre que si $\overline X$ est de t.f.
\deleted{quasi-proj}
sur $\overline{\mathbf Q}$, l'hom $\mathcal G_{\overline X}
\mathrel{\tikz[baseline=-.45ex,x=1em,y=1em]{
  \draw (0,0)--(1.5,0);
  \fill (1.5,0)--(1.36,.08)--(1.36,0)--cycle;
}}\Gamma$ a une image
d'indice fini (puisqu'on peut
\tikz[baseline=0pt]{
  \node[inner sep=0pt] at (.12,.60) {\inserted{alors}};
  \draw (-.35,-.13) .. controls (-.30,.25) and (.28,.31) .. (.48,.39);
}
réduire le
corps de base à une ext. finie $K$ de $\mathbf Q$).

\textcircled{2} \uline{Supposons $X$ connexe
\tikz[baseline=0pt]{
  \node[inner sep=0pt] at (.18,.53) {\inserted{$\ne\varnothing$}};
  \draw (-.23,-.18) .. controls (-.35,.06) and (.55,.27) .. (.78,.80);
}, et}
Considérons un revêtement universel $\widetilde X\xrightarrow{\unsure{h|b}}X$
$X$, et \deleted{considérons} posons
\[
(5)\qquad \mathcal E(X,\widetilde X)\ \text{ou}\ \mathcal E_X=
\operatorname{Aut}\left(\widetilde X
% Literal overloaded separator; deletion state and interpretation remain doubtful.
\unsure{\mathbin{\tikz[baseline=-.45ex,x=1em,y=1em]{
  \draw[line width=.65pt] (.06,-.32)--(.26,.34);
  \draw[line width=.85pt] (-.03,0)--(.42,0);
  \fill (.42,0)--(.31,.065)--(.31,-.065)--cycle;
}}}X\right)
\]
\mspage{1}
